\subsection{Objects, reduction modulo $p$ and the motivic group}
\label{ssec:mumfordobject}

\Cref{thm:mumfordnumerical} reduces the semiregularity route to one object at
one CM point. We first derive what any object with the right Chern character
must satisfy, whether or not it is semiregular: no sequence of standard
operations produces it from line bundles; it extends along no first order
deformation of $D^{b}(Y)$ other than the direction of the curve; if it is
semiregular, it splits only into summands with no $\Ext^{2}$ between them;
and its self-extensions are large in every degree, with an odd part that the
Euler characteristic forces beyond what Hochschild cohomology supplies. We
then turn to reduction modulo $p$. By Li's proof of the Tate conjecture for
powers of abelian fourfolds over finite fields, the classes are represented
by cycles at every closed point of every reduction of the curve, and the
target is equivalent to a bound on those cycles on a Zariski dense set of
closed points. Finally, the motivic group of $X$ is one of four groups, and
the kinds of algebraic classes known on products of $X$ with other abelian
varieties generate no exceptional class. None of this constructs such an
object or proves the target.

\begin{lemma}[The contraction ranks are symmetric]\label{lem:rankduality}
Let $A$ be an abelian variety of dimension $g$ and $c\in H^{*}(A,\CC)$ a class
each of whose components has Hodge type $(p,p)$ for some $p$. For $k\ge0$ put
$r_{k}(c)=\dim\bigl(HT^{k}(A)\lrcorner\,c\bigr)$, with the action of
\Cref{setup:hh}. Then $r_{k}(c)=r_{g-k}(c)$ for $0\le k\le g$, and
$r_{k}(c)=0$ for $k>g$.
\end{lemma}

\begin{proof}
Write $H^{1,0}=H^{1,0}(A)$ and $A'=H^{0,1}(A)$, both of dimension $g$, so
that $H^{*}(A,\CC)=\bigwedge^{*}H^{1,0}\otimes\bigwedge^{*}A'$ and
$HT^{*}(A)=\bigwedge^{*}(A'\oplus T)$ with $T=(H^{1,0})^{\vee}$. Fix a
generator of $\bigwedge^{g}A'$ and let
$\star\colon\bigwedge^{q}A'\to\bigwedge^{g-q}A'^{\vee}$ send $\eta$ to
$\zeta\mapsto\eta\wedge\zeta$. Under $\Phi=\mathrm{id}\otimes\star$, wedging
with $z\in A'$ becomes, up to a sign fixed on each bigraded piece, contraction
with $z$ on the second factor, while $T$ still contracts the first. So $\Phi$
identifies the action of $HT^{*}(A)$ with the action of
$\bigwedge^{*}U^{\vee}$ by interior products on $\bigwedge^{*}U$, where
$U=H^{1,0}\oplus A'^{\vee}$ has dimension $2g$, up to signs that do not change
ranks. A class of type $(p,p)$ goes to
$\bigwedge^{p}H^{1,0}\otimes\bigwedge^{g-p}A'^{\vee}\subset\bigwedge^{g}U$,
so $\eta=\Phi(c)$ lies in the middle degree whatever the $p$ of its
components. Hence $r_{k}(c)$ is the rank of
$\bigwedge^{k}U^{\vee}\to\bigwedge^{g-k}U$, $x\mapsto\iota_{x}\eta$, which is
zero for $k>g$. Its transpose for the natural pairings is
$\bigwedge^{g-k}U^{\vee}\to\bigwedge^{k}U$, $y\mapsto\pm\iota_{y}\eta$,
because $\langle\iota_{x}\eta,y\rangle=\langle\eta,x\wedge y\rangle
=\pm\langle\iota_{y}\eta,x\rangle$. A map and its transpose have the same
rank.
\end{proof}

\begin{theorem}[The self-extensions of an object with an exceptional Chern
character]\label{thm:mumfordprofile}
Let $c\in C$ and $Y=X_{c}\times X_{c}$, and write
$e_{k}=\dim\Ext^{k}(E,E)$ for a perfect complex $E$ on $Y$.
\begin{enumerate}[label=\textup{(\roman*)},leftmargin=2.6em,itemsep=2pt]
\item For every $E$ and every $k$, $e_{k}\ge r_{k}(\ch E)$ and
  $e_{k}=e_{8-k}$.
\item For a rational exceptional class $\omega=\omega_{a}$ with
  $L_{0}(a)\ne0$: $r_{0}=r_{8}=1$, $r_{1}=r_{7}=16$ and $r_{2}=r_{6}=119$.
\item Let $\omega$ be a rational exceptional class and $\ch(E)=N\omega$
  with $N\ne0$. Then
  \[
    e_{1}+e_{3}=e_{0}+e_{2}+\tfrac12e_{4}
      +\sum_{j\ge1}(-1)^{j}e_{-j}.
  \]
  If $\Ext^{<0}(E,E)=0$, for instance if $E$ is a sheaf, then
  $\sum_{k\ \mathrm{odd}}e_{k}=\sum_{k\ \mathrm{even}}e_{k}$, and
  $e_{1}+e_{3}\ge120+\frac12e_{4}$ when $L_{0}(a)\ne0$.
\end{enumerate}
\end{theorem}

\begin{proof}
(i) The identity $\sigma_{E}\circ\operatorname{ev}_{E}=(\,\cdot\,)\lrcorner
\ch(E)$ of \Cref{setup:hh} holds in every degree, so the kernel of
$\operatorname{ev}_{E}$ on $HT^{k}(Y)$ lies in the annihilator of $\ch(E)$,
and the image of $HT^{k}(Y)$ in $\Ext^{k}(E,E)$ has dimension at least
$r_{k}(\ch E)$. Serre duality with the trivial canonical bundle gives
$e_{k}=e_{8-k}$.

(ii) $r_{0}=1$ because $\omega\ne0$, and $r_{2}=119$ is
\Cref{thm:mumfordrigid}(ii). If $w\in HT^{1}(Y)$ annihilates $\omega$, so does
$w'\wedge w$ for every $w'\in HT^{1}(Y)$, since the annihilator is a left
ideal; for $w\ne0$ these products span a space of dimension $15$, which cannot
lie in $\Ann_{HT^{2}}(\omega)=\CC\kappa$. So $r_{1}=16$, and
\Cref{lem:rankduality} gives $r_{8},r_{7},r_{6}$.

(iii) Since $\ch(E)$ has only a component in $H^{4}$, $\ch(E^{\vee})\ch(E)$
lies in $H^{8}(Y)$ and $\chi(E,E)=\int_{Y}\ch(E^{\vee})\ch(E)=0$, the Todd
class being trivial. By (i) and Serre duality, which also pairs
$\Ext^{-j}$ with $\Ext^{8+j}$,
\[
  0=\chi(E,E)=2(e_{0}-e_{1}+e_{2}-e_{3})+e_{4}
  +2\sum_{j\ge1}(-1)^{j}e_{-j},
\]
which is the displayed identity. With no negative extensions the odd and even
parts are equal, and $e_{0}\ge1$, $e_{2}\ge119$ by (i) and (ii) when
$L_{0}(a)\ne0$.
\end{proof}

\begin{proposition}[Divisor classes do not reach the exceptional classes]
\label{prop:mumforddivisor}
Let $t\in C$ and let $\omega$ be a rational exceptional class. Then $\omega_{t}$
is not in the $\QQ$-span of products of divisor classes of $X_{t}\times
X_{t}$: that span meets the rational Hodge classes of $H^{2}\otimes H^{2}$
exactly in the classes with $\alpha\in\QQ$. Consequently no perfect complex
whose Chern character lies in the subring generated by divisor classes, in
particular no bounded complex of direct sums of line bundles, has
$\ch(E)=N\omega$ with $N\ne0$. At a CM point $c$ the fourfold $X_{c}$ has eight
Hodge classes in $H^{4}$, of which the products of divisor classes span six.
\end{proposition}

\begin{proof}
At a point with Hodge group $G$, the divisor classes of $X_{t}\times X_{t}$ are
the three classes of \Cref{prop:mumfordproduct}(ii), all invariant under
$\Sp(V,\psi)$, and so are their products; the $\Sp(V,\psi)$-invariant part of
$H^{2}\otimes H^{2}$ is spanned by $\pi_{0}$ and $\pi_{12}+\pi_{13}+\pi_{23}$,
the classes with $a_{1}=a_{2}=a_{3}$. At a CM point the Hodge group is a
maximal torus of $G$ (proof of \Cref{thm:mumfordcm}), so the Hodge classes
are the monomials of weight zero. The proof of \Cref{prop:cmsource}(ii)
counts them: $16$ divisor classes $e^{(i)}_{w}e^{(j)}_{-w}$, $132$ Hodge
classes in degree four, of which the products of divisor classes span $100$,
and modulo that span
$\omega_{a}\equiv\frac12\bigl((a_{2}-a_{3})u_{12}+(a_{1}-a_{3})u_{13}
+(a_{1}-a_{2})u_{23}\bigr)$ with $u_{12},u_{13},u_{23}$ linearly independent.
So the span meets that of the four $\pi$ in exactly the classes with
$a_{1}=a_{2}=a_{3}$. For a rational class these
coordinates are the conjugates of $\alpha$, which by \Cref{lem:conjugates} are
equal only if $\alpha\in\QQ$. The Chern character of a complex in the subring
generated by divisor classes has its component in $H^{4}$ in the span of
products of two divisor classes, which excludes $N\omega$. For $X_{c}$ itself
the monomials of weight zero in $H^{4}$ are the six products
$e_{w}e_{-w}e_{w'}e_{-w'}$ over two opposite pairs, which are products of
divisor classes, and the two tetrahedron monomials.
\end{proof}

At a CM point the Hodge classes of $X_{c}$ in $H^{4}$ that are not products
of divisor classes are among the classes that the proof of
\Cref{thm:mumfordcm} pulls back along homomorphisms $X_{c}\times X_{c}\to
X_{c}$; they are algebraic by Markman's theorem \cite{Mar25a}. Li proves,
independently and more generally, that every Hodge class on every power of
every complex CM abelian fourfold is algebraic \cite[Theorem~1.3]{Li26}, which
gives \Cref{thm:mumfordcm} for every power $X_{c}^{n}$ and not only for the
square.

By \Cref{prop:mumforddivisor}, an object with $\ch(E)=N\omega$ cannot have
its Chern character in the subring generated by divisor classes, so it cannot
be assembled from line bundles alone. The next theorem shows that no sequence
of the standard operations of the derived category, applied to line bundles
on powers of $X_{t}$, gets any further.

\begin{theorem}[Objects generated by line bundles]\label{thm:lefschetzclosure}
Let $t\in C$, $X=X_{t}$, and let $L(X)\subseteq\Sp(V,\psi)$ be the Lefschetz
group of $X$, the centraliser of $\End^{0}(X)$ in $\Sp(V,\psi)$. Let
$\mathcal{A}$ be the class of abelian varieties isogenous to a power of $X$;
through an isogeny $X^{n}\to B$ the group $L(X)$ acts diagonally on
$H^{*}(B,\QQ)$, and the action does not depend on the isogeny. Let
$\mathcal{L}$ be the smallest class of objects of the categories
$D^{b}(B)$, $B\in\mathcal{A}$, that contains the line bundles and is closed
under shifts, cones, derived tensor products, derived duals, translations,
the functors $Lf^{*}$ and $Rf_{*}$ for homomorphisms $f$ between members of
$\mathcal{A}$, and the Fourier--Mukai functors $\Phi_{K}\colon D^{b}(B)\to
D^{b}(B')$ with kernel $K\in\mathcal{L}$ on $B\times B'$.
\begin{enumerate}[label=\textup{(\roman*)},leftmargin=2.6em,itemsep=2pt]
\item Every object of $\mathcal{L}$ has an $L(X)$-invariant Chern character.
\item $L(X)=\Sp(V,\psi)$ if $t$ is not a CM point, and $L(X)$ is a maximal
  torus of $\Sp(V,\psi)$ if it is. In both cases the $L(X)$-invariant classes
  on $X\times X$ form the subring generated by divisor classes, of dimensions
  $1,3,6,10,15,10,6,3,1$, respectively $1,16,100,304,454,304,100,16,1$, in
  the degrees $0,2,\dots,16$.
\item No rational exceptional class $\omega_{t}$ is $L(X)$-invariant. Hence
  for no object $E\in\mathcal{L}$ on $X\times X$ and no $N\ne0$ is
  $\ch(E)-N\omega_{t}$ invariant under $L(X)$; in particular
  $\ch(E)\ne N\omega_{t}$.
\item The $\Sp(V,\psi)$-module generated by $\pi_{12}-\pi_{13}$ is
  irreducible, of highest weight $2\varpi_{2}$ and dimension $308$, and it
  meets the span of the four $\pi$ in the plane of the classes with
  $a_{0}=0$ and $a_{1}+a_{2}+a_{3}=0$.
\end{enumerate}
\end{theorem}

\begin{proof}
(i) Two isogenies $X^{n}\to B$ differ by an element of $\End^{0}(X^{n})=
M_{n}(\End^{0}(X))$, which commutes with $L(X)$ acting diagonally, so the
action is well defined. For the same reason a homomorphism $f$ between members
of $\mathcal{A}$ is $L(X)$-equivariant on $H^{1}$, so $f^{*}$ is equivariant
on $H^{*}$, and so is $f_{*}$, which is $f^{*}$ conjugated by Poincar\'e
duality, since $L(X)\subseteq\SL(V)$ fixes the orientation classes.
Translations act trivially on cohomology. The Chern character is additive on
distinguished triangles and multiplicative on derived tensor products,
$\ch(E^{\vee})=\sum_{i}(-1)^{i}\ch_{i}(E)$, $\ch(Lf^{*}E)=f^{*}\ch(E)$, and
$\ch(Rf_{*}E)=f_{*}\ch(E)$ by the Grothendieck--Riemann--Roch theorem
\cite[Theorem~15.2]{Ful98}, the Todd class of an abelian variety being $1$;
and $\Phi_{K}=Rp_{2*}(p_{1}^{*}(\,\cdot\,)\otimes^{L}K)$. So the objects with
$L(X)$-invariant Chern character form a class with all the closure
properties, and it remains to see that it contains the line bundles, that is,
that $c_{1}(\mathcal{O}(D))$ is $L(X)$-invariant. Pulled back along an
isogeny, which is injective on rational cohomology, $c_{1}$ becomes a Hodge
class of degree two on some $X^{n}$, an element of the invariants of the
Hodge group in ${\textstyle\bigwedge^{2}}(V^{\vee}\otimes\QQ^{n})=
{\textstyle\bigwedge^{2}}V^{\vee}\otimes S^{2}\QQ^{n}\oplus S^{2}V^{\vee}
\otimes{\textstyle\bigwedge^{2}}\QQ^{n}$. If $t$ is not a CM point, the Hodge
group is $G$, whose invariants in ${\textstyle\bigwedge^{2}}V$ are the multiples of $\psi$
and in $S^{2}V$ are zero, since every summand of $S^{2}(V_{1}\otimes V_{2}
\otimes V_{3})$ contains a factor $S^{2}V_{i}$; so the Hodge classes of degree
two are $\Sp(V,\psi)$-invariant. If $t$ is a CM point, a monomial of degree two
in a weight basis has weight zero for the Hodge torus exactly when its two
weights are opposite, and then it has weight zero for $L(X)$, by the
description in (ii).

(ii) If $t$ is not a CM point, $\End^{0}(X)=\QQ$ and $L(X)=\Sp(V,\psi)$. At a
CM point $\End^{0}(X)\otimes\CC$ is the algebra of diagonal matrices in a
weight basis (proof of \Cref{thm:mumfordcm}), and $\psi$ pairs each weight
with its opposite only, so $L(X)_{\CC}$ is the torus of diagonal matrices
$\operatorname{diag}(\lambda_{w})$ with $\lambda_{w}\lambda_{-w}=1$, of rank
four. For $t$ not CM, the first fundamental theorem for $\Sp(V,\psi)$
\cite[Chapter~5]{GW09} shows that its invariants in
${\textstyle\bigwedge^{*}}(V\otimes\QQ^{2})$ are spanned by the products of
the three quadratic invariants, which are the divisor classes. Their
dimensions follow from the skew Cauchy formula
${\textstyle\bigwedge^{2m}}(V\otimes\QQ^{2})=\bigoplus_{\lambda}S_{\lambda}V
\otimes S_{\lambda'}\QQ^{2}$, over the partitions $\lambda=(2^{r},1^{s})$ of
$2m$ with at most two columns, and the rule that $S_{\lambda}V$ contains the
trivial representation of $\Sp(V,\psi)$, then once, exactly when both columns
of $\lambda$ have even length at most $8$ \cite{GW09}: the invariants of
degree $2m$ have dimension $\sum(s+1)$, over the even $r,s\ge0$ with
$2r+s=2m$ and $r+s\le8$, which is $1,3,6,10,15,10,6,3,1$ for
$m=0,\dots,8$. At a CM point the invariants are spanned by the
monomials of weight zero for $L(X)$; a set of weights $\pm e_{j}$ with sum
zero is a union of opposite pairs, so each such monomial is, up to sign, a
product of monomials of degree two and weight zero, which are divisor
classes. For each opposite pair $\{w,-w\}$ the four vectors
$e^{(1)}_{\pm w}$, $e^{(2)}_{\pm w}$ give one monomial of weight zero in
degree $0$, four in degree $2$ and one in degree $4$, so the numbers of these
monomials are the coefficients of $(1+4y+y^{2})^{4}$.

(iii) If $t$ is not a CM point, the $\Sp(V,\psi)$-invariant combinations of the
$\pi$ are those with $a_{1}=a_{2}=a_{3}$ (\Cref{prop:mumforddivisor}), and
the coefficients of $\omega_{t}$ are the distinct conjugates of $\alpha$
(\Cref{lem:conjugates}). At a CM point use the map $x\mapsto E_{x}$ of the
proof of \Cref{lem:oneexceptional}, which is $\Sp(V,\psi)$-equivariant: it
carries $\omega_{t}$ to an element $(r,\beta)$ of $\QQ\times F$ with
$\beta\notin\QQ$, acting on $U_{ij}$ by the conjugate $\sigma_{k}(\beta)$
(\Cref{prop:mumfordrm}(ii)), and $\omega_{t}$ is $L(X)$-invariant only if
$E_{\omega_{t}}$ preserves the weight spaces of $L(X)$ in ${\textstyle\bigwedge^{2}}V$. Take
weight vectors $v=x\otimes y\otimes z$ and $w=x\otimes y'\otimes z'$ of
weights $(1,1,1)$ and $(1,-1,-1)$ for the Hodge torus. They are not
opposite, so $v\wedge w$ spans a weight space of $L(X)$, every weight of
$L(X)$ on $V$ having multiplicity one. But
$v\wedge w=(x\cdot x)\otimes\bigl(y\cdot y'\otimes z\wedge z'+y\wedge y'
\otimes z\cdot z'\bigr)$ has nonzero components in $U_{12}$ and in $U_{13}$,
on which $E_{\omega_{t}}$ acts by $\sigma_{3}(\beta)\ne\sigma_{2}(\beta)$. So
the line is not preserved. The rest of (iii) follows from (i).

(iv) Write $H^{2}(X)=\QQ\theta\oplus P$, with $P$ the primitive part, the
irreducible $\Sp(V,\psi)$-module of highest weight $\varpi_{2}$ and dimension
$27$. Through the form ${\textstyle\bigwedge^{2}}\psi$ the four $\pi$ are
self-adjoint endomorphisms of $H^{2}(X)$, and $\pi_{12}-\pi_{13}$ kills
$\theta$ and is a trace-free self-adjoint endomorphism of $P$, an element of
$S^{2}P$ orthogonal to the invariant. Restricting
$S^{2}({\textstyle\bigwedge^{2}}V)=S_{(2,2)}V\oplus{\textstyle\bigwedge^{4}}V$
from $\GL(V)$ to $\Sp(V,\psi)$ gives, with dimensions $336=308+27+1$ and
$70=42+27+1$, the decomposition $S^{2}P=V(2\varpi_{2})\oplus V(\varpi_{4})
\oplus V(\varpi_{2})\oplus\QQ$. The element $\pi_{12}-\pi_{13}$ is
$G$-invariant. The summand $V(\varpi_{2})\cong P$ has no $G$-invariant, since
$({\textstyle\bigwedge^{2}}V)^{G}=\QQ\theta$, and neither has $V(\varpi_{4})$,
which lies in ${\textstyle\bigwedge^{4}}V$, whose $G$-invariants are the line
of $\theta^{2}$ (\Cref{prop:mumfordproduct}(i)), invariant under all of
$\Sp(V,\psi)$. So $\pi_{12}-\pi_{13}$ is a nonzero vector of the irreducible
module $V(2\varpi_{2})$, which it generates; its dimension is $308$ by Weyl's
dimension formula. The same argument puts every trace-free combination of
$\pi_{12},\pi_{13},\pi_{23}$ in $V(2\varpi_{2})$, while $\pi_{0}$ and
$\pi_{12}+\pi_{13}+\pi_{23}$ are $\Sp(V,\psi)$-invariant and independent
modulo it; so adding the four $\pi$ raises the dimension by two.
\end{proof}

\Cref{fig:hodgecount} compares the Hodge classes of $X_{t}\times X_{t}$ with
the subring generated by divisor classes, degree by degree.

\begin{figure}[htbp]
\centering
\includegraphics[width=\textwidth]{fig_hodgecount}
\caption{\Cref{thm:lefschetzclosure}(ii). The Hodge classes of
$X_{t}\times X_{t}$ degree by degree, against the subring $D^{p}$ generated
by divisor classes, which contains the Chern character of every object of
$\mathcal{L}$: at a very general point of the Mumford curve (left) and at a
CM point (right). The hatched classes, the exceptional ones among them, are
reached by no object of $\mathcal{L}$.}
\label{fig:hodgecount}
\end{figure}

\begin{remark}[Where a construction has to leave $\mathcal{L}$]
\label{rem:leavel}
\Cref{thm:lefschetzclosure} extends \Cref{prop:mumforddivisor} from sums of
line bundles to everything the standard operations make out of them, at every
point of the curve, including the CM points, where the exceptional classes
are algebraic. There the cycles that represent them are not Chern characters
of objects of $\mathcal{L}$, and neither are the two Hodge classes of $X_{c}$
in degree four that are not products of divisor classes. At a point that is
not CM, every Chern character of an object of $\mathcal{L}$ lies in the
isotypic component of the trivial representation of $\Sp(V,\psi)$, while by
\Cref{thm:lefschetzclosure}(iv) an object with $\ch(E)=N\omega$ must have a
nonzero component in the one of highest weight $2\varpi_{2}$. Every
construction of such an object therefore contains a step that does not commute with the Lefschetz group on
cohomology: passing to a direct summand, a subobject or a Jordan--H\"older
factor, choosing a point of a moduli space, or deforming from a member with
more endomorphisms. The last is what semiregularity does along the curve
(\Cref{thm:mumfordnumerical}).
\end{remark}

\begin{lemma}[The summands of a semiregular object]\label{lem:semiregsum}
Let $Z$ be a smooth projective variety and $E_{1},E_{2}$ perfect complexes on
$Z$. On $\Ext^{2}(E_{1}\oplus E_{2},E_{1}\oplus E_{2})=\bigoplus_{i,j}
\Ext^{2}(E_{j},E_{i})$ the semiregularity map is
$\sigma_{E_{1}\oplus E_{2}}(\zeta)=\sigma_{E_{1}}(\zeta_{11})
+\sigma_{E_{2}}(\zeta_{22})$. So $\sigma_{E_{1}\oplus E_{2}}$ is injective if
and only if $\sigma_{E_{1}}$ and $\sigma_{E_{2}}$ are injective, their images
meet only in zero, and $\Ext^{2}(E_{1},E_{2})=\Ext^{2}(E_{2},E_{1})=0$.
\end{lemma}

\begin{proof}
The Atiyah class is natural, so $\mathrm{At}(E_{1}\oplus E_{2})$ is the
diagonal sum of $\mathrm{At}(E_{1})$ and $\mathrm{At}(E_{2})$. In
$\sigma(\zeta)=\sum_{q}\frac1{q!}\operatorname{tr}(\zeta\circ\mathrm{At}^{q})$
of \Cref{thm:BF} the composite $\zeta\circ\mathrm{At}^{q}$ has diagonal blocks
$\zeta_{ii}\circ\mathrm{At}(E_{i})^{q}$, and the trace of a block matrix of
morphisms is the sum of the traces of its diagonal blocks. In particular the
off-diagonal blocks $\zeta_{12}$ and $\zeta_{21}$ lie in the kernel, so
injectivity forces them to vanish; on the diagonal blocks the map is
$(\zeta_{11},\zeta_{22})\mapsto\sigma_{E_{1}}(\zeta_{11})+\sigma_{E_{2}}
(\zeta_{22})$, which is injective exactly when both summands are injective
and their images meet only in zero.
\end{proof}

The condition on the images cannot be dropped, and the objects that show it
are the simplest ones: the structure sheaves of abelian subvarieties of half
dimension, each of which is semiregular on its own.

\begin{lemma}[Abelian subvarieties of half dimension are semiregular]
\label{lem:absemireg}
Let $A$ be an abelian variety of dimension $2n$ and $B\subseteq A$ an abelian
subvariety of dimension $n$. Then $\operatorname{ev}_{\cO_{B}}\colon
HT^{2}(A)\to\Ext^{2}(\cO_{B},\cO_{B})$ is surjective, its kernel is
$\Ann_{HT^{2}}([B])$, of dimension $6n^{2}-n$, and $\sigma_{\cO_{B}}$ is
injective.
\end{lemma}

\begin{proof}
The normal bundle of $B$ is trivial, so $\ch(\cO_{B})=[B]$ by the
Grothendieck--Riemann--Roch theorem for the embedding, and
$\Ext^{*}(\cO_{B},\cO_{B})\cong{\textstyle\bigwedge^{*}}\CC^{2n}$ (proof of
\Cref{cor:naturalsums}), so $\dim\Ext^{2}(\cO_{B},\cO_{B})=\binom{2n}{2}
=2n^{2}-n$, while $\dim HT^{2}(A)=2\binom{2n}{2}+4n^{2}=8n^{2}-2n$. Choose a
complementary abelian subvariety $B'$; the isogeny $B\times B'\to A$
identifies $H^{1}(A,\CC)$ with $H^{1}(B)\oplus H^{1}(B')$ compatibly with the
Hodge decompositions, identifies $HT^{2}(A)$ with $HT^{2}(B\times B')$
compatibly with contraction, and carries $[B]$ to a nonzero multiple of
$1\otimes[\mathrm{pt}_{B'}]$, where $[\mathrm{pt}_{B'}]$ spans
${\textstyle\bigwedge^{2n}}H^{1}(B')$. Write $T=T_{B}\oplus T_{B'}$ for the
tangent space and $H^{0,1}(A)=H^{0,1}(B)\oplus H^{0,1}(B')$. A class
$\zeta\in H^{0,2}(A)={\textstyle\bigwedge^{2}}H^{0,1}(A)$ has
$\zeta\cup[B]=\zeta_{BB}\otimes[\mathrm{pt}_{B'}]$, with $\zeta_{BB}$ its
component in ${\textstyle\bigwedge^{2}}H^{0,1}(B)$, because every factor from
$H^{0,1}(B')$ is killed by the top class of $B'$; so the annihilator in
$H^{0,2}$ has dimension $\binom{2n}{2}-\binom{n}{2}$. For $a\otimes v\in
H^{0,1}(A)\otimes T=H^{1}(T_{A})$ the contraction is $a\wedge(v\lrcorner[B])
=a_{B}\otimes(v_{B'}\lrcorner[\mathrm{pt}_{B'}])$, since $v_{B}$ kills
$[\mathrm{pt}_{B'}]$ and $a_{B'}\wedge(v_{B'}\lrcorner[\mathrm{pt}_{B'}])$
would have type $(n-1,n+1)$ on the $n$-fold $B'$; the map
$v\mapsto v_{B'}\lrcorner[\mathrm{pt}_{B'}]$ has rank $n$, as does
$a\mapsto a_{B}$, so the contraction has rank $n^{2}$ on the space
$H^{0,1}(A)\otimes T$ of dimension $4n^{2}$ and the annihilator there has
dimension $3n^{2}$. For $\vartheta\in{\textstyle\bigwedge^{2}}T$ only the
component in ${\textstyle\bigwedge^{2}}T_{B'}$ contracts $[\mathrm{pt}_{B'}]$
nontrivially, and injectively; so the annihilator in
${\textstyle\bigwedge^{2}}T$ has dimension $\binom{2n}{2}-\binom{n}{2}$.
Altogether $\dim\Ann_{HT^{2}}([B])=2\binom{2n}{2}-2\binom{n}{2}+3n^{2}
=6n^{2}-n$. By the square $\sigma_{E}\circ\operatorname{ev}_{E}=
(\,\cdot\,)\lrcorner\ch(E)$ of \Cref{setup:criterion} the kernel of
$\operatorname{ev}_{\cO_{B}}$ lies in $\Ann_{HT^{2}}([B])$, so
$\operatorname{ev}_{\cO_{B}}$ has rank at least $(8n^{2}-2n)-(6n^{2}-n)
=2n^{2}-n=\dim\Ext^{2}(\cO_{B},\cO_{B})$: it is surjective and its kernel is
the annihilator. If $\sigma_{\cO_{B}}(y)=0$, write $y=\operatorname{ev}(x)$;
then $x\lrcorner[B]=0$, so $x$ lies in the annihilator, which is the kernel,
and $y=0$.
\end{proof}

\begin{remark}[Sums of abelian subvarieties]\label{rem:absums}
Let $B_{1},\dots,B_{r}\subseteq A$ be abelian subvarieties of dimension
$n\ge3$, translated so that any two of them meet in finitely many points or
not at all, which generic translates achieve, and let $E=\bigoplus_{j}
\cO_{B_{j}}[k_{j}]$. The sheaves $\mathcal{E}xt^{q}(\cO_{B_{i}},\cO_{B_{j}})$
for $i\ne j$ are supported on $B_{i}\cap B_{j}$ and vanish for $q\ne n$,
by the Koszul resolution of $\cO_{B_{i}}$ restricted to $B_{j}$, so
$\Ext^{2}(\cO_{B_{i}}[k_{i}],\cO_{B_{j}}[k_{j}])=H^{2+k_{j}-k_{i}-n}
(B_{i}\cap B_{j},\mathcal{E}xt^{n})$ vanishes when all $k_{j}$ are equal and
$n\ge3$, and when $|k_{i}-k_{j}|\le1$ and $n\ge4$; and each
$\sigma_{\cO_{B_{j}}}$ is injective by \Cref{lem:absemireg}. In these cases
every hypothesis of \Cref{lem:semiregsum}, applied to one summand and the
sum of the others, holds except the one on the images. Now let $v\in HT^{2}(A)$
satisfy $v\lrcorner\ch(E)=\sum_{j}(-1)^{k_{j}}v\lrcorner[B_{j}]=0$ with
$v\lrcorner[B_{1}]\ne0$. Then $\zeta=\bigoplus_{j}\operatorname{ev}_{\cO_{B_{j}}}(v)$
is nonzero, because $v$ is not in the annihilator of $[B_{1}]$, and
$\sigma_{E}(\zeta)=\sum_{j}(-1)^{k_{j}}v\lrcorner[B_{j}]=0$: the images of
the $\sigma_{\cO_{B_{j}}}$ meet, and $E$ is not semiregular. This is exactly
the situation at a member of a Weil family on which a Weil class is a
combination of classes of abelian subvarieties of dimension $n$, as on a
product of elliptic curves: a tangent direction $v$ of the family keeps
$\ch(E)$ of Hodge type but moves any $[B_{j}]$ outside
$\QQ\theta^{n}\oplus\HW(A)$, and \Cref{cor:naturalsums} excludes such sums
for this reason. The sum is obstructed although its summands are not.
\end{remark}

\begin{corollary}[Hochschild deformations and direct summands]
\label{cor:mumfordsplit}
Let $c\in C$, $Y=X_{c}\times X_{c}$, $\omega=\omega_{a}$ a rational
exceptional class with $L_{0}(a)\ne0$, and $E$ a perfect complex on $Y$ with
$\ch(E)=N\omega$, $N\ne0$.
\begin{enumerate}[label=\textup{(\roman*)},leftmargin=2.6em,itemsep=2pt]
\item $E$ extends to the first order deformation of $D^{b}(Y)$ given by
  $x\in HH^{2}(Y)=HT^{2}(Y)$ only if $x\in\CC\kappa$. No $B$-field, no
  Poisson bivector and no commutative deformation off the curve carries it.
\item If $\sigma_{E}$ is injective and $E=E_{1}\oplus E_{2}$ with both
  summands nonzero, then $\Ext^{2}(E_{1},E_{2})=\Ext^{2}(E_{2},E_{1})=0$, each
  $\sigma_{E_{i}}$ is injective, each $E_{i}$ extends to first order along the
  curve, and $\kappa\lrcorner\ch(E_{i})=0$. If moreover
  $\dim\Ext^{2}(E,E)=119$, then $\dim\Ext^{2}(E_{i},E_{i})=
  120-\dim\Ann_{HT^{2}}(\ch E_{i})$, and
  \[
    \Ann_{HT^{2}}(\ch E_{1})+\Ann_{HT^{2}}(\ch E_{2})=HT^{2}(Y),\qquad
    \Ann_{HT^{2}}(\ch E_{1})\cap\Ann_{HT^{2}}(\ch E_{2})=\CC\kappa .
  \]
\end{enumerate}
\end{corollary}

\begin{proof}
(i) The kernel of $\operatorname{ev}_{E}$ on $HH^{2}(Y)$ is the space of
deformations of $D^{b}(Y)$ along which $E$ deforms to first order
\cite{Tod09}, \cite[Section~8.3]{Mar25a}. If $\operatorname{ev}_{E}(x)=0$ then
$N\,x\lrcorner\,\omega=\sigma_{E}(\operatorname{ev}_{E}(x))=0$, so
$x\in\CC\kappa$ by \Cref{thm:mumfordrigid}(ii). The three kinds of
deformation named are the three summands of $HT^{2}(Y)$.

(ii) The first two assertions are \Cref{lem:semiregsum}. By
\Cref{thm:mumfordnumerical}(ii), $\operatorname{ev}_{E}(\kappa)=0$, and
$\operatorname{ev}_{E}$ is the diagonal sum of $\operatorname{ev}_{E_{1}}$
and $\operatorname{ev}_{E_{2}}$, so each vanishes at $\kappa$; then
$\kappa\lrcorner\ch(E_{i})=\sigma_{E_{i}}(\operatorname{ev}_{E_{i}}(\kappa))
=0$. Put $e(x)=120-\dim\Ann_{HT^{2}}(x)$, the rank of contraction into $x$.
It is subadditive in $x$, and by \Cref{thm:mumfordprofile}(i)
$\dim\Ext^{2}(E_{i},E_{i})\ge e(\ch E_{i})$. The cross terms vanish, so
\[
  119=\dim\Ext^{2}(E,E)=\sum_{i}\dim\Ext^{2}(E_{i},E_{i})
  \ge e(\ch E_{1})+e(\ch E_{2})\ge e(N\omega)=119,
\]
and every inequality is an equality. Hence
$\dim\Ann(\ch E_{1})+\dim\Ann(\ch E_{2})=121$. The intersection of the two
annihilators lies in $\Ann(N\omega)=\CC\kappa$ and contains $\kappa$, so it is
$\CC\kappa$, and the sum has dimension $120$.
\end{proof}

Part (ii) plays for the Mumford classes the role of \Cref{thm:nosum}: a
semiregular object attaining the bound either does not split, or splits into
two pieces whose annihilators fill $HT^{2}(Y)$ and meet only in the direction
of the curve. The two annihilators are the analogue of the two branches $P$
and $Q$ of \Cref{rem:twobranches}.

\begin{proposition}[Tate classes at every closed point of every reduction]
\label{prop:modp}
There are a finitely generated subring $R\subset\CC$, a model $W_{R}\to
C_{R}$ of $W\to C$ over $R$, with $C_{R}$ smooth and projective over $R$ with
geometrically connected fibres and $W_{R}$ an abelian scheme, projective over
$C_{R}$, and a finite surjective morphism $C'_{R}\to C_{R}$ from a second curve of the same
kind, with the following properties. Put $Y'=(W_{R}\times_{C_{R}}W_{R})
\times_{C_{R}}C'_{R}$, with projection $\pi'\colon Y'\to C'_{R}$, and let
$\omega$ be a rational exceptional class. For every prime $\ell$ there is a
global section $\omega_{\ell}$ of the lisse sheaf
$R^{4}\pi'_{*}\QQ_{\ell}(2)$ on $C'_{R}[1/\ell]$ whose restriction to
$C'_{R}\otimes_{R}\CC$ is the pull-back of $\omega\otimes1$; write
$\omega_{\ell}(\bar{x})$ for its value at a geometric point over a point $x$.
\begin{enumerate}[label=\textup{(\roman*)},leftmargin=2.6em,itemsep=2pt]
\item At a closed point $x$ of $C'_{R}[1/\ell]$ the residue field
  $\kappa(x)$ is finite, $Y'_{x}$ is the square of an abelian fourfold over
  it, and $\omega_{\ell}(\bar{x})$ is fixed by $\Gal(\bar{\kappa}/\kappa(x))$.
\item At such a point $\omega_{\ell}(\bar{x})$ is a $\QQ_{\ell}$-linear
  combination of classes of algebraic cycles of codimension two on $Y'_{x}$,
  defined over $\kappa(x)$.
\item Let $s$ be a closed point of $\Spec R[1/\ell]$ and $\eta_{s}$ the generic
  point of the reduction $C'_{s}$. Then $\omega_{\ell}(\bar\eta_{s})$ is fixed
  by the Galois group of the function field $\kappa(\eta_{s})$, and it is not
  a $\QQ_{\ell}$-linear combination of products of divisor classes of
  $Y'_{\bar\eta_{s}}$.
\end{enumerate}
Every closed point $x$ of a reduction $C_{s}$ of $C$ is the image of a closed
point $x'$ of $C'_{s}$, and $Y'_{x'}=Y_{x}\otimes_{\kappa(x)}\kappa(x')$; so
(i) and (ii) hold at every closed point of every reduction of $C$, after a
finite extension of its residue field.
\end{proposition}

\begin{proof}
\emph{The models.} The family $W\to C$ is defined over a subfield
$k_{1}\subset\CC$ finitely generated over $\QQ$, say as $W_{1}\to C_{1}$. Let
$K_{1}$ be the function field of $C_{1}$ and fix an embedding over $k_{1}$ of
an algebraic closure $\bar{K}_{1}$ into $\CC$; it sends the geometric generic
point of $C_{1}$ to a point $t\in C$, and $Y_{t}=X_{t}\times X_{t}$ is the base
change of $Y_{\bar{K}_{1}}$. The class $\omega_{t}$ is a Hodge class on an
abelian variety, hence absolutely Hodge, so it comes from an absolutely Hodge
class of $Y_{\bar{K}_{1}}$, and $\Gal(\bar{K}_{1}/K_{1})$ acts on these
through a finite quotient \cite[Theorem~2.11 and Proposition~2.9]{Del82}. Let
$K'\subset\bar{K}_{1}$ be a finite extension of $K_{1}$ on which the action is
trivial, $k'$ the algebraic closure of $k_{1}$ in $K'$, and $C'$ the smooth
projective curve over $k'$ with function field $K'$; it is geometrically
connected over $k'$ and maps finitely onto $C_{1}\otimes_{k_{1}}k'$.
Spreading out $W_{1}\to C_{1}$, $C'$ and this map over a finitely generated
subring of $k'$, and inverting one element of it, gives a normal $R$ and
models with the stated properties.

\emph{The section.} By smooth and proper base change
$R^{4}\pi'_{*}\QQ_{\ell}(2)$ is lisse on $C'_{R}[1/\ell]$. Its stalk at the
geometric generic point $\Spec\bar{K}_{1}$ is $H^{4}(Y_{\bar{K}_{1}},
\QQ_{\ell}(2))=H^{4}(Y_{t},\QQ)\otimes\QQ_{\ell}$, and $\omega_{t}\otimes1$
is the $\ell$-adic component of an absolutely Hodge class, so it is fixed by
$\Gal(\bar{K}_{1}/K')$. The scheme $C'_{R}[1/\ell]$ is normal and integral
with function field $K'$, so $\Gal(\bar{K}_{1}/K')$ maps onto its fundamental
group \cite[Tag~0BQI]{Stacks}, and the vector extends to a global section
$\omega_{\ell}$. On the connected curve $C'_{R}\otimes_{R}\CC$ it agrees with
the pull-back of $\omega\otimes1$, the two being global sections of the same
local system that agree at a point over $t$.

(i) and (ii). A scheme of finite type over $\ZZ$ has finite residue fields at
its closed points. A global section of a lisse sheaf is fixed by the
fundamental group of every point, which is (i); and (ii) is Li's theorem
\cite[Theorem~1.1]{Li26}, with $n=r=2$, for the abelian fourfold that is the
fibre of $W_{R}$ at the image of $x$, base changed to $\kappa(x)$.

(iii) The first assertion holds for the same reason as (i). For the second,
let $\bar\eta$ be a geometric generic point of $C'_{R}$ and $\bar{s}$ a
geometric point over $s$. The morphism $C'_{R}\to\Spec R$ is flat and proper
with geometrically connected and geometrically reduced fibres, so the
specialisation map $\pi_{1}(C'_{\bar\eta})\to\pi_{1}(C'_{\bar{s}})$ is
surjective \cite[Tag~0C0P]{Stacks}, and it is compatible with the actions of
the two groups on the first cohomology of the fibres of $W_{R}$, identified by
cospecialisation. The image of $\pi_{1}(C'_{\bar\eta})$ contains that of
$\pi_{1}(C'_{\CC})$, the closure of the image of the topological fundamental
group of a finite cover of $C$, whose Zariski closure contains
$G_{\QQ_{\ell}}$ because the connected monodromy group of $W\to C$ is $G$
(proof of \Cref{cor:mumfordlefschetz}). So the Zariski closure of the
monodromy of $W_{R}$ over $C'_{\bar{s}}$ contains $G_{\QQ_{\ell}}$. A divisor
class of $Y'_{\bar\eta_{s}}$ is defined over a finite extension of
$\kappa(\eta_{s})$, so it is fixed by an open subgroup of the Galois group,
whose image in $\pi_{1}(C'_{s})$ is open \cite[Tag~0BQI]{Stacks} and meets
$\pi_{1}(C'_{\bar{s}})$ in an open subgroup; the Zariski closure of the image
of an open subgroup has the same identity component, so the class is
$G_{\QQ_{\ell}}$-invariant. The $G$-invariants of ${\textstyle\bigwedge^{2}}(V\oplus V)$ are the
three divisor classes of \Cref{prop:mumfordproduct}(ii), which are
$\Sp(V,\psi)$-invariant, so every product of divisor classes of
$Y'_{\bar\eta_{s}}$ is $\Sp(V,\psi)$-invariant. The class
$\omega_{\ell}(\bar\eta_{s})$ corresponds to $\omega\otimes1$, whose
coefficients $a_{1},a_{2},a_{3}$ are the distinct conjugates of $\alpha$
(\Cref{lem:conjugates}), and the $\Sp(V,\psi)$-invariant classes among the
$\omega_{a}$ are those with $a_{1}=a_{2}=a_{3}$ (\Cref{prop:mumforddivisor}).
The last sentence of the statement holds because $C'_{R}\to C_{R}$ is finite
and surjective.
\end{proof}

At the generic point of each reduction the class is therefore a Tate class
over a global function field that is not a combination of products of divisor
classes, exactly as over $\CC$, while at every closed point it is a
combination of classes of cycles. The next theorem shows that a bound on those
cycles is all that is missing, with no passage through $p$-adic Hodge theory.

\begin{theorem}[The Mumford target is a degree bound in positive
characteristic]\label{thm:modpdensity}
Keep \Cref{prop:modp}, fix a prime $\ell$ and a relatively ample line bundle
on $Y'\to C'_{R}$, and for a finite set $\Phi$ of tuples of Hilbert
polynomials say that $\omega$ is \emph{$\Phi$-represented} at a point $x$ of
$C'_{R}[1/\ell]$ if $\omega_{\ell}(\bar{x})=\sum_{i}q_{i}\operatorname{cl}(Z_{i})$
for some $q_{i}\in\QQ_{\ell}$ and closed subschemes $Z_{i}\subset Y'_{\bar{x}}$
of pure codimension two whose Hilbert polynomials form a tuple in $\Phi$. The
following are equivalent.
\begin{enumerate}[label=\textup{(\alph*)},leftmargin=2.6em,itemsep=2pt]
\item $\omega_{t}$ is algebraic for every $t\in C$.
\item For some finite $\Phi$, the closed points of $C'_{R}[1/\ell]$ at which
  $\omega$ is $\Phi$-represented are Zariski dense in $C'_{R}$.
\item For some finite $\Phi$ and some Zariski dense set $S$ of closed points of
  $\Spec R[1/\ell]$, the class $\omega$ is $\Phi$-represented at infinitely
  many closed points of $C'_{s}$ for every $s\in S$.
\end{enumerate}
For a single closed point $s$ of $\Spec R[1/\ell]$, the class $\omega$ is
$\Phi$-represented at infinitely many closed points of $C'_{s}$, for some
finite $\Phi$, if and only if $\omega_{\ell}(\bar\eta_{s})$ is a
$\QQ_{\ell}$-linear combination of classes of algebraic cycles on
$Y'_{\bar\eta_{s}}$.
\end{theorem}

\begin{proof}
One standard property of the $\ell$-adic cycle class is used beyond those over
a field: if $\mathcal{O}$ is a discrete valuation ring in which $\ell$ is
invertible, $\mathcal{Y}\to\Spec\mathcal{O}$ is smooth and projective, and
$\mathcal{Z}\subset\mathcal{Y}$ is a closed subscheme flat over $\mathcal{O}$,
then the specialisation isomorphism of $\ell$-adic cohomology carries the
cycle class of the geometric generic fibre of $\mathcal{Z}$ to that of its
geometric special fibre \cite[Cycle]{SGA4h}. The specialisation isomorphism
also carries the values of a global section of a lisse sheaf to each other,
by its construction.

\emph{Step 1: the represented loci are closed.} For a tuple
$\underline{P}=(P_{1},\dots,P_{m})$ let $\cH_{\underline{P}}$ be the fibre
product over $C'_{R}[1/\ell]$ of the relative Hilbert schemes
$\Hilb^{P_{i}}(Y'/C'_{R}[1/\ell])$. It is projective over $C'_{R}[1/\ell]$
\cite{Gro61}, so it has finitely many connected components. On a component
$T$, with universal subschemes $\mathcal{Z}_{i}$, put $A(h)=\{q\in
\QQ_{\ell}^{m}:\sum_{i}q_{i}\operatorname{cl}(\mathcal{Z}_{i,\bar{h}})
=\omega_{\ell}(\bar{h})\}$ for $h\in T$. If $h'$ is a specialisation of $h$,
there is a discrete valuation ring mapping to $T$ with generic point to $h$ and
closed point to $h'$ \cite[Tag~054F]{Stacks}, and the property above, applied
to the pull-backs of $Y'$ and of the $\mathcal{Z}_{i}$, gives $A(h)=A(h')$.
Any two points of the connected noetherian scheme $T$ are joined by a chain of
specialisations and generisations, through the generic points of its
irreducible components, so $A$ is constant on $T$, say $A_{T}$. Let
$\Sigma_{\underline{P}}$ be the union of the images in $C'_{R}[1/\ell]$ of the
components $T$ with $A_{T}\ne\emptyset$. It is closed, as a finite union of
images of proper morphisms, and it is the set of points at which $\omega$ is
$\{\underline{P}\}$-represented: at a point $x$ of the image of such a $T$, a
closed point of the fibre of $T$ over $x$ has residue field finite over
$\kappa(x)$ and gives subschemes of $Y'_{\bar{x}}$. Put
$\Sigma_{\Phi}=\bigcup_{\underline{P}\in\Phi}\Sigma_{\underline{P}}$.

\emph{Step 2: (b) implies (a).} $\Sigma_{\Phi}$ is closed and contains a dense
set, so it is $C'_{R}[1/\ell]$ and contains the generic point $\xi$ of
$C'_{R}$, whose residue field is $K'$. So $\omega_{t}\otimes1=
\omega_{\ell}(\bar\xi)=\sum_{i}q_{i}\operatorname{cl}(Z_{i})$ with
$q\in\QQ_{\ell}^{m}$ and subschemes $Z_{i}$ of $Y_{\bar{K}_{1}}$, and so of
$Y_{t}$. The cycle class maps are compatible with the comparison isomorphism
$H^{4}(Y_{t},\QQ)\otimes\QQ_{\ell}\cong H^{4}(Y_{t},\QQ_{\ell})$
\cite[Section~1 and Example~2.1(a)]{Del82}, so the linear system
$\omega_{t}=\sum_{i}q_{i}[Z_{i}]$, which has rational coefficients, is
solvable over $\QQ_{\ell}$, hence over $\QQ$, and $\omega_{t}$ is algebraic.
The $Z_{i}$ are defined over a finite extension of $K'$, so they spread out to
subschemes, flat over a nonempty open subset \cite[Tag~052A]{Stacks}, of the
pull-back of $W\times_{C}W$ to a finite cover of $C$, and by the local
constancy of the cycle class used in the proof of \Cref{thm:closedlocus},
$\omega$ is algebraic at every point of a nonempty Zariski open subset of $C$,
which is uncountable. \Cref{cor:mumfordequiv} gives (a).

\emph{Step 3: (c) implies (b).} Let $D$ be the set of closed points at which
$\omega$ is $\Phi$-represented and $\bar{D}$ its closure. For $s\in S$ the set
$D\cap C'_{s}$ is infinite, so its closure is the irreducible curve $C'_{s}$,
and $C'_{s}\subseteq\bar{D}$. The points $u$ of $\Spec R[1/\ell]$ over which
the fibre of $\bar{D}$ has dimension at least one form a closed set, by the
upper semicontinuity of the dimension of the fibres of a proper morphism
\cite[Tag~0D4I]{Stacks}; it contains the dense set $S$, so it contains the
generic point. The generic fibre of $\bar{D}$ is then a closed subset of
dimension one of the irreducible curve $C'_{R}\otimes_{R}k'$, so it contains
$\xi$, and $\bar{D}=C'_{R}[1/\ell]$.

\emph{Step 4: (a) implies (c).} By (a) and Step 1, applied over $\bar{K}_{1}$
with the Betti class in place of the $\ell$-adic one, $\omega_{t}$ is the
class of a rational combination of subschemes of $Y_{\bar{K}_{1}}$, defined
over a finite extension $L$ of $K'$. Enlarge $R$ to a finitely generated
subring of the algebraic closure of $k'$ in $L$, let $C''_{R}\to C'_{R}$ be the
normalisation in $L$, and spread the subschemes out over $C''_{R}$; by generic
flatness \cite[Tag~052A]{Stacks} they are flat over a dense open subset $U$ of
$C''_{R}[1/\ell]$, with constant Hilbert polynomials. By Step 1 applied on $U$,
$\omega$ is represented by their fibres at every point of $U$. The image of $U$
in $C'_{R}[1/\ell]$ contains a dense open subset $U'$, every closed point of
$U'$ is the image of a closed point of $U$, and $U'$ meets $C'_{s}$ in a
nonempty open, hence infinite, subset for every $s$ in a dense open subset of
$\Spec R[1/\ell]$.

\emph{The single reduction.} If $\omega$ is $\Phi$-represented at infinitely
many closed points of $C'_{s}$, then $\Sigma_{\Phi}\cap C'_{s}$ is closed and
infinite, so it is $C'_{s}$ and contains $\eta_{s}$. Conversely, cycles at
$\bar\eta_{s}$ are defined over a finite extension of $\kappa(\eta_{s})$, and
spreading them out as in Step 4 over a finite cover of $C'_{s}$ represents
$\omega$, with fixed Hilbert polynomials, at every closed point of a nonempty
open subset of $C'_{s}$.
\end{proof}

\begin{remark}[What is left in positive characteristic]\label{rem:modpremains}
By \Cref{prop:modp}(ii) the class is represented at every closed point of
$C'_{R}[1/\ell]$, so \Cref{thm:modpdensity} turns the Mumford target into a
statement of uniformity alone: the Hilbert data of some representation must
stay in one finite set on a Zariski dense set of closed points, for instance
at infinitely many closed points of infinitely many reductions. The passage
to characteristic zero is then done by the specialisation of cycle classes
along $C'_{R}$, and neither the $p$-adic variational Hodge conjecture nor a
lifting of cycles is needed. Li's theorem gives no bound: where its proof
carries cycles into a prescribed isogeny class, it does so by an isogeny whose
degree it does not control \cite[proof of Theorem~1.2, Step~3]{Li26}.

Nor do the natural representatives at the supersingular points supply the
bound. At a closed point $x$ whose fibre is supersingular, the Frobenius
eigenvalues on $H^{2}(Y'_{\bar{x}},\QQ_{\ell}(1))$ are roots of unity, so by
Tate's theorem \cite{Tat66} this space is spanned by divisor classes, and
$\omega_{\ell}(\bar{x})\in{\textstyle\bigwedge^{4}}H^{1}$ is a combination of
products of two of them. But for a closed point $s$ of $\Spec R[1/\ell]$ and
a finite set of Hilbert polynomials, only finitely many closed points of
$C'_{s}$ carry such a combination $\sum_{i}q_{i}\operatorname{cl}(D_{i})
\operatorname{cl}(D'_{i})$ with effective divisors $D_{i},D'_{i}$ whose
Hilbert polynomials lie in that set: otherwise Step~1 of the proof of
\Cref{thm:modpdensity}, run on relative Hilbert schemes of divisors, would
give one at $\bar\eta_{s}$, against \Cref{prop:modp}(iii). This holds whether
or not the supersingular points are dense. They are dense in $C'_{R}$ when $R$
has dimension one. Indeed, for every prime $p\ge5$ not dividing the
discriminants of $F$ and $D$, which is all but finitely many, the
supersingular points of the reduction at a prime over $p$ of Kisin's integral
model of the Shimura curve, with hyperspecial level at $p$, are its whole
Newton jumping locus and number $(p^{r}-1)(g-1)$ on each component of genus
$g$, $r$ being the local degree of $F$ \cite[Theorem~1.2 and
Corollary~5.17]{SZ11}. For all but finitely many $s$ the curve $C'_{s}$ maps
finitely onto such a component, of genus at least two, so it has at least
$p-1$ supersingular points, while a closed subset of $C'_{R}$ of dimension one
meets all but finitely many of the $C'_{s}$ in boundedly many points.

Within one reduction the bound is equivalent to the class being algebraic at
the generic point of $C'_{s}$ (\Cref{thm:modpdensity}), a Tate class over a
global function field that is not a combination of products of divisor
classes (\Cref{prop:modp}(iii)). This is the codimension two case of the
variational Tate conjecture along $C'_{s}$, which Morrow proves for divisors
in its crystalline form \cite{Mor19}. From one reduction alone, the way back
to characteristic zero would be the $p$-adic variational Hodge conjecture for
$W\times_{C}W$ over the Witt vectors of $\kappa(s)$, with rational rather than
$\ell$-adic coefficients and with the crystalline class of the cycle matching
the de Rham class of $\omega$. Bloch, Esnault and Kerz lift such a class in
$K_{0}\otimes\QQ$ of the special fibre to the formal completion for $p>15$,
the dimension being nine, and the algebraisation of the formal class is the
remaining step of their programme \cite[Theorem~1.3]{BEK14}. The route
therefore runs through the uniform bound or through this conjecture.
\end{remark}

Five of the routes of \Cref{rem:bypassaudit} survive the rigidity theorems:
\emph{specialisation}, \emph{reduction modulo $p$}, \emph{semiregularity},
\emph{Kuga--Satake} and \emph{motivated classes}. The next corollary shows
that they are not five separate obligations.

\begin{corollary}[The five routes reduce to one statement]\label{cor:mumfordone}
Let $W\to C$ be as in \Cref{cor:mumfordlefschetz}, let $\omega$ be a rational
exceptional class, and let the \emph{Mumford target} be the first case of
\textup{(F2)}: the Hodge conjecture for $X_{t}\times X_{t}$ for every $t\in C$.
\begin{enumerate}[label=\textup{(\roman*)},leftmargin=2.6em,itemsep=2pt]
\item Each of the following is equivalent to the Mumford target: $\omega_{t}$ is
  algebraic for uncountably many $t\in C$; $B(W\times_{C}W)$ holds for
  $L_{W\times_{C}W}$; $\omega$ is represented with Hilbert data in one finite
  set at infinitely many points of $C$, for instance at infinitely many CM
  points; $\omega$ is represented with Hilbert data in one finite set on a
  Zariski dense set of closed points of the arithmetic model of
  \Cref{thm:modpdensity}.
\item Each of the following implies the Mumford target: a perfect complex $E$
  at one point of $C$ with $\ch(E)=N\omega$, $N\ne0$, no other component, and
  $\dim\Ext^{2}(E,E)=119$; the algebraicity, for uncountably many $t$, of the
  Kuga--Satake class of the surface $S_{\lambda}$ of \Cref{thm:mumfordks}
  attached to $X_{t}$; a rational map as in \Cref{prop:hksquare}(iv), under
  its hypotheses, from $X_{t}\times X_{t}$ for uncountably many $t$.
\item There is a countable set $\Sigma_{0}\subset C$ containing the CM points
  such that the Mumford target holds as soon as the Hodge conjecture holds for
  $X_{c}\times X_{c}$ at one point $c\notin\Sigma_{0}$.
\end{enumerate}
\end{corollary}

\begin{proof}
(i) The first two are \Cref{cor:mumfordequiv}, with \Cref{lem:oneexceptional}
for one class in place of two, the third is \Cref{cor:mumfordbounded} and the
fourth is \Cref{thm:modpdensity}. (ii) The first is
\Cref{thm:mumfordnumerical}(iii). The second and the third give the Hodge
conjecture for $X_{t}\times X_{t}$ at each of the uncountably many points, by
\Cref{thm:mumfordks}(ii) and \Cref{prop:hksquare}(iv), and then (i) applies.
(iii) In the proof of \Cref{cor:mumfordbounded} the set of points at which
$\omega$ is algebraic is the union of the countably many closed sets
$\Sigma_{\underline{P},\underline{q}}$, each finite or equal to $C$. Let
$\Sigma_{0}$ be the union of the finite ones and of the CM points. If
$\omega_{c}$ is algebraic at a point $c\notin\Sigma_{0}$, then $c$ lies in
some $\Sigma_{\underline{P},\underline{q}}$ that is not finite, which is $C$,
and (i) applies.
\end{proof}

\begin{remark}[Which routes are needed]\label{rem:mumfordneeded}
Of the five routes, specialisation, reduction modulo $p$ and motivated
classes are the Mumford target itself in three other forms, by
\Cref{cor:mumfordone}(i) and, for the motivated classes,
\Cref{cor:mumfordequiv}; semiregularity and Kuga--Satake are stronger inputs
that imply it. So any one route suffices. Each of the three
forms is the whole of the problem, a bound on Hilbert data or a case of
\textup{(L)} that is equivalent to the target, and the two stronger inputs
ask for an object and for a cycle.
\end{remark}

The target can also be read at a single member, through the motivic group of
$X$ rather than through cycles on the family; it then becomes the question of
which of four groups this is.

\begin{proposition}[Four possible motivic groups]\label{prop:mumfordmotivic}
In \Cref{setup:mumford}, let $M$ be the group of the Tannakian category of
homological motives generated by $h(X)$, with the Betti realisation as fibre
functor, let $M^{1}=M\cap\Sp(V,\psi)$, and let $N$ be the normaliser of $G$ in
$\Sp(V,\psi)$. Over $\bar{\QQ}$ identify $V_{1}$, $V_{2}$ and $V_{3}$ with one
plane, carrying $\epsilon_{1}$, $\epsilon_{2}$, $\epsilon_{3}$ to one form, and
let $\zeta\in\Sp(V,\psi)(\bar{\QQ})$ be the cyclic permutation
$v_{1}\otimes v_{2}\otimes v_{3}\mapsto v_{3}\otimes v_{1}\otimes v_{2}$ of
the three factors. It normalises $G$, and other identifications change it by
an element of $G(\bar{\QQ})$; so it acts through $V$ on the Hodge classes of
every power of $X$, independently of the choice.
\begin{enumerate}[label=\textup{(\roman*)},leftmargin=2.6em,itemsep=2pt]
\item $N/G$ is a form of the symmetric group $S_{3}$,
  $M=\mathbb{G}_{m}\cdot M^{1}$, and $M^{1}$ is one of $G$, $G\cdot A_{3}$,
  $N$ and $\Sp(V,\psi)$, where $G\cdot A_{3}$ is the inverse image in $N$ of
  the subgroup of order three of $N/G$. The last three contain $\zeta$.
\item The following are equivalent: $M^{1}=G$; the Hodge conjecture holds for
  $X^{n}$ for every $n\ge1$; the two exceptional classes of $X\times X$ are
  algebraic; some algebraic class on some power of $X$ is not fixed by
  $\zeta$.
\item Let $\mathrm{Det}\in(V^{\otimes4})^{G}\subset H^{4}(X^{4},\QQ)$ be the
  Hodge class of \Cref{rem:mumforddet}, and write $r_{1}$ also for the Hodge
  class $t\in(V^{\otimes4})^{G}$ with $\hat{t}=r_{1}=s_{1}+s_{2}+s_{3}$, in the
  notation of Steps~2 and~3 of the proof of \Cref{thm:mumfordpowers}. Both
  are fixed by $N$ and not by $\Sp(V,\psi)$, and each of them is algebraic if
  and only if $M^{1}\ne\Sp(V,\psi)$. So the algebraicity of either excludes
  the last case only; by this argument it does not give the exceptional
  classes, which are not algebraic when $M^{1}$ is $G\cdot A_{3}$ or $N$.
\item Let $W\to C$ be as in \Cref{cor:mumfordlefschetz}, and suppose that it
  is defined over a number field; this holds for the restriction of the
  universal family of a fine moduli space of polarised abelian fourfolds with
  level structure to the normalisation of one of the Shimura curves of
  \cite{Mum69} in it, and for its pull-backs to finite \'etale covers. Then the
  set $\Sigma_{0}$ constructed in the proof of \Cref{cor:mumfordone}(iii) lies
  in $C(\bar{\QQ})$, and the Mumford target holds as soon as, at one point
  $t\in C(\CC)\setminus C(\bar{\QQ})$, some algebraic cycle on some power of
  $X_{t}$ has a class that is not fixed by $\zeta$.
\end{enumerate}
\end{proposition}

\begin{proof}
\emph{The group.} On an abelian variety homological and numerical
equivalence coincide \cite{Lie68}, so the homological motives generated by
$h(X)$ and the Lefschetz motive form a semisimple abelian category
\cite{Jan92}. The K\"unneth projectors are algebraic \cite[\S2]{Kle68}, so
after the usual change of sign in the commutativity constraint this is a
Tannakian category over $\QQ$, with the Betti realisation as a fibre functor,
faithful by the definition of homological equivalence. Its group $M$ is
reductive, the category being semisimple, and the $M$-invariants in the
realisation of a motive $P$ are the image of $\Hom(\mathbf{1},P)$
\cite[Theorem~2.11]{DM82}; so for all $m$ and $p$ the $M$-invariants of
$H^{2p}(X^{m},\QQ)(p)$ are the classes of algebraic cycles. The group $M$ acts
faithfully on $V\oplus\QQ(1)$ and fixes the polarisation class, so it lies in
$\mathrm{GSp}(V,\psi)$ and acts on $\QQ(1)$ through the multiplier. Through
$\psi$ and the K\"unneth projectors every tensor space of $V\oplus\QQ(1)$ is a
sum of direct summands of spaces $H^{k}(X^{k},\QQ)(j)$ cut out by algebraic
projectors, so the $M$-invariants there are algebraic classes, hence Hodge
classes, fixed by $\MT(X)=\mathbb{G}_{m}\cdot G$; as a reductive group is the
subgroup fixing its invariants in the tensor spaces
\cite[Proposition~3.1]{Del82}, $\MT(X)\subseteq M$ and $G\subseteq M^{1}$. An
element of $M$ with multiplier $\nu$ is a square root of $\nu$ times an
element of $M^{1}$, so $M=\mathbb{G}_{m}\cdot M^{1}$. The homotheties act
trivially on the spaces $H^{2p}(X^{m},\QQ)(p)$, so there the algebraic classes
are the $M^{1}$-invariants, as the Hodge classes are the $G$-invariants.

\emph{The groups between $G$ and $\Sp(V,\psi)$.} Over $\bar{\QQ}$ the Lie
algebra $\mathfrak{sp}(V,\psi)\cong S^{2}V$ is the sum of
$\operatorname{Lie}G=T_{1}\oplus T_{2}\oplus T_{3}$, where $T_{i}\cong
S^{2}V_{i}\otimes{\textstyle\bigwedge^{2}}V_{j}\otimes{\textstyle\bigwedge^{2}}V_{k}$, and of
$S^{2}V_{1}\otimes S^{2}V_{2}\otimes S^{2}V_{3}$, which is an irreducible
representation of $G$ of dimension $27$; so a $G$-stable subspace of
$\mathfrak{sp}(V,\psi)$ containing $\operatorname{Lie}G$ is
$\operatorname{Lie}G$ or $\mathfrak{sp}(V,\psi)$. If $H$ is a closed
subgroup with $G\subseteq H\subseteq\Sp(V,\psi)$, its Lie algebra is such a
subspace. If it is $\operatorname{Lie}G$, then $H^{\circ}=G$, which $H$
normalises, so $H\subseteq N$; if it is $\mathfrak{sp}(V,\psi)$, then
$H^{\circ}=\Sp(V,\psi)$, which is connected, so $H=\Sp(V,\psi)$.

\emph{The normaliser.} Since $V$ is absolutely irreducible, the centraliser of
$G$ in $\Sp(V,\psi)$ is $\{\pm1\}$, which lies in $G$ as the image of
$(-1,1,1)$, and $G(\bar{\QQ})$ maps onto the group of inner automorphisms of
$G_{\bar{\QQ}}$. So $N(\bar{\QQ})/G(\bar{\QQ})$ embeds in the outer
automorphism group of $G_{\bar{\QQ}}$, which permutes the three factors,
$\SL_{2}$ having no outer automorphism. The six permutations of the factors,
after identifications as above, preserve $\psi=\epsilon_{1}\otimes
\epsilon_{2}\otimes\epsilon_{3}$ and normalise $G$, so
$N(\bar{\QQ})/G(\bar{\QQ})\cong S_{3}$. The Galois group acts on it by
conjugation through its action on the three factors, which is transitive, and
the subgroups of $S_{3}$ stable under conjugation by a transitive subgroup are
$1$, $A_{3}$ and $S_{3}$. So the algebraic subgroups of $N$ containing $G$ are
$G$, $G\cdot A_{3}$ and $N$, and with the previous paragraph this proves (i),
$\zeta$ mapping to an element of order three of $N/G$.

(ii) If $M^{1}=G$, then $M=\MT(X)$, and every Hodge class on every power of
$X$ is algebraic. The Hodge conjecture for $X\times X$ includes the
exceptional classes. If these are algebraic, so is every rational exceptional
class $\omega_{a}$ of \Cref{setup:omegaa}, and $\omega_{a}$ is not fixed by
$\zeta$: $\zeta$ permutes the three factors cyclically, so it fixes $\theta$
and carries $U_{ij}=S^{2}V_{i}\otimes S^{2}V_{j}\otimes{\textstyle\bigwedge^{2}}V_{k}$
onto $U_{\zeta(i)\zeta(j)}$; labelling the factors so that $\zeta$ sends the
$i$-th to the $(i+1)$-th, it carries $\pi_{12}$ to $\pi_{23}$, $\pi_{23}$ to
$\pi_{13}$ and $\pi_{13}$ to $\pi_{12}$ and fixes $\pi_{0}$. So it fixes
$\omega_{a}$ only if $a_{1}=a_{2}=a_{3}$, whereas these are the distinct
conjugates of an irrational element of $F$ (\Cref{lem:conjugates}). Finally,
if $M^{1}\ne G$, then $M^{1}$ contains $\zeta$ by (i), and every algebraic
class, being $M^{1}$-invariant, is fixed by $\zeta$.

(iii) The class $\mathrm{Det}$ spans $(S^{4}V)^{G}\otimes
{\textstyle\bigwedge^{4}}\QQ^{4}$, where $(S^{4}V)^{G}$ is a line
(\Cref{rem:mumforddet}). As $N$ normalises $G$, it preserves this line, and it
fixes it, because $N(\bar{\QQ})$ is generated by $G(\bar{\QQ})$ and the
permutations of the factors, which fix Cayley's hyperdeterminant, a
symmetric function of the three factors, and the restriction of
$\mathrm{Det}$ to the diagonal is a nonzero multiple of it
(\Cref{rem:mumforddet}). $\Sp(V,\psi)$
fixes no nonzero element of $S^{4}V$, which would be a nonzero quartic form
constant on the single orbit $V\setminus\{0\}$. The permutations of the
factors conjugate $s_{k}$ into $s_{\sigma(k)}$, so they commute with $r_{1}$,
as does $G$; and $t\mapsto\hat{t}$ is $\Sp(V,\psi)$-equivariant, since the
pairing $\int_{X}v\cup w\cup\theta^{3}$ on $V$ is $G$-invariant and
alternating, hence a multiple of the form $\psi$. So the class $r_{1}$ is
fixed by $N$. On $S^{2}V\subset V\otimes V$ the endomorphism $r_{1}$ acts by
$3$ on $S^{2}V_{1}\otimes S^{2}V_{2}\otimes S^{2}V_{3}$ and by $-1$ on the
$T_{i}$, since $s_{k}$ is $1$ on $S^{2}V_{k}$ and $-1$ on
${\textstyle\bigwedge^{2}}V_{k}$; as $S^{2}V$ is irreducible under
$\Sp(V,\psi)$, the class $r_{1}$ is not $\Sp(V,\psi)$-invariant. If
$M^{1}\ne\Sp(V,\psi)$, then $M^{1}\subseteq N$ by (i), and the two classes are
$M^{1}$-invariant, hence algebraic; if $M^{1}=\Sp(V,\psi)$, they are not. When
$M^{1}$ is $G\cdot A_{3}$ or $N$ it contains $\zeta$, which fixes no
exceptional class.

(iv) A fine moduli space as in the statement is defined over a number field,
and the curve contains infinitely many CM points, which are algebraic, an
abelian variety with complex multiplication being defined over a number
field; so the curve is the Zariski closure of a set of algebraic points and
is defined over $\bar{\QQ}$, and so are its normalisation, the restricted
family and their finite \'etale covers. Now let $W\to C$ be defined over
$\bar{\QQ}$. The moduli map of the family is a finite morphism defined over
$\bar{\QQ}$, so the CM points of $C$, whose images are algebraic, lie in
$C(\bar{\QQ})$. The other points of $\Sigma_{0}$ lie in the finite sets
$\Sigma_{\underline{P},\underline{q}}$ of the proof of
\Cref{cor:mumfordbounded}, images of unions of connected components of fibre
products of relative Hilbert schemes of $W\times_{C}W\to C$. These are base
changes of schemes over $\bar{\QQ}$, and each connected component of the
base change is the base change of a connected component over $\bar{\QQ}$,
which is algebraically closed; so a finite
$\Sigma_{\underline{P},\underline{q}}$ consists of algebraic points. If
$t\notin C(\bar{\QQ})$, then $t\notin\Sigma_{0}$, the Hodge group of $X_{t}$ is
$G$, and (ii) for $X_{t}$ gives the Hodge conjecture for $X_{t}\times X_{t}$;
\Cref{cor:mumfordone}(iii) gives the target.
\end{proof}

The four cases differ already on $X^{4}$ and $X^{6}$
(\Cref{fig:mumfordgroups}). Over $\bar{\QQ}$,
$(V^{\otimes2m})^{G}$ is the tensor product of three copies of the space of
$\SL_{2}$-invariants of $(\bar{\QQ}^{2})^{\otimes2m}$, of dimension $c=2$ for
$m=2$ and $c=5$ for $m=3$, with a basis of non-crossing pairings that $\zeta$
and the transpositions of the factors permute. So the algebraic classes in
the multilinear part $V^{\otimes2m}\subset H^{2m}(X^{2m},\QQ)$ have dimension
$c^{3}$, $(c^{3}+2c)/3$, $(c^{3}+3c^{2}+2c)/6$ and $(2m-1)(2m-3)\cdots1$ when
$M^{1}$ is $G$, $G\cdot A_{3}$, $N$ and $\Sp(V,\psi)$, the last by the first
fundamental theorem for $\Sp(V,\psi)$: $8$, $4$, $4$, $3$ on $X^{4}$ and
$125$, $45$, $35$, $15$ on $X^{6}$. The middle two count the orbits of the
basis under $\langle\zeta\rangle$ and under $S_{3}$: $\zeta$ fixes the $c$
triples of equal pairings and a transposition fixes $c^{2}$ of them. If $M$ were known to be
connected, $M^{1}$ would be connected too, since $-1\in G$ lies in its
identity component; it would then be $G$ or $\Sp(V,\psi)$, and the
algebraicity of $\mathrm{Det}$ alone would give the Hodge conjecture for every
power of $X$. The argument does not decide whether $M$ is connected. It is the
standard use of the motivic group as the stabiliser of the algebraic classes,
and we claim no priority for \Cref{prop:mumfordmotivic}. It proves no new case
of the Hodge conjecture: it says what one new cycle on one power of one member
would have to do.

\begin{figure}[tbp]
\centering
  \includegraphics[width=\textwidth]{fig_mumfordgroups}
\caption{The four possible groups $M^{1}=M\cap\Sp(V,\psi)$ of
\Cref{prop:mumfordmotivic}, where $M=\mathbb{G}_{m}\cdot M^{1}$, as the chain
$G\subset G\cdot A_{3}\subset N\subset\Sp(V,\psi)$: the first two inclusions
have index $3$ and $2$, and $N$, whose identity component is $G$, has
dimension $9$ against $36$. For each case the table gives the dimension of
the algebraic classes in the multilinear part $V^{\otimes2m}$ of
$H^{2m}(X^{2m},\QQ)$ for $m=2$ and $m=3$, from the $c=2$, respectively $c=5$,
non-crossing pairings; whether $M^{1}$
contains the cyclic permutation $\zeta$ of the three tensor factors; whether
$\mathrm{Det}$ and $r_{1}$ are algebraic; and whether the two exceptional
classes of $X\times X$ are. Only $M^{1}=G$ makes the exceptional classes
algebraic, and it is equivalent to the Hodge conjecture for every power of
$X$ (\Cref{prop:mumfordmotivic}(ii)); the proposition does not decide which
of the four cases occurs. Below, from left to right: $\zeta$; the subgroups of
$N(\bar{\QQ})/G(\bar{\QQ})\cong S_{3}$, of which only $1$, $A_{3}$ and $S_{3}$
are stable under the Galois group, which acts transitively on the three
factors; and $\mathfrak{sp}(V,\psi)=\operatorname{Lie}G\oplus
S^{2}V_{1}\otimes S^{2}V_{2}\otimes S^{2}V_{3}$, $36=9+27$, the second summand
irreducible under $G$, with the eigenvalues $-1$ and $3$ of $r_{1}$ on the two
summands.}
\label{fig:mumfordgroups}
\end{figure}

The same invariance explains why the classes known or expected to be
algebraic on abelian varieties stop short of the exceptional classes, on all
abelian varieties at once.

\begin{proposition}[The system generated by the known classes]\label{prop:mumfordformal}
Keep \Cref{setup:mumford}, and $\zeta$ and $G\cdot A_{3}$ as in
\Cref{prop:mumfordmotivic}. For an abelian variety $A$ let $A_{X}\subseteq A$
be the largest abelian subvariety isogenous to a power of $X$, and $A'$ the
identity component of the intersection of the kernels of the homomorphisms
$A\to X$. Then $A_{X}\times A'\to A$ is an isogeny, $\Hom(X,A')=0$, and
\[
  H^{1}(A,\QQ)=\bigl(V\otimes\Hom^{0}(A_{X},X)\bigr)\oplus H^{1}(A',\QQ),
\]
the first summand embedded by $v\otimes f\mapsto f^{*}v$. Let $\GL(V)$ act on
$H^{*}(A,\QQ)={\textstyle\bigwedge^{*}}H^{1}(A,\QQ)$ through the factor $V$. Let $\cS$
be the smallest family of subspaces $\cS(A)\subseteq H^{*}(A,\QQ)$, over all
complex abelian varieties $A$, that is closed under linear combinations, cup
products, exterior products, and pull-back and push-forward along
homomorphisms, and contains:
\begin{enumerate}[label=\textup{(\alph*)},leftmargin=2.6em,itemsep=1pt]
\item the divisor classes;
\item for every subfield $E\subseteq\End^{0}(A)$, the rational classes whose
  extension to $\bar{\QQ}$ lies in $\bigoplus_{\sigma}{\textstyle\bigwedge^{d}}
  H^{1}(A,\bar{\QQ})_{\sigma}$, the sum over the embeddings
  $\sigma\colon E\to\bar{\QQ}$ of the top exterior powers of the eigenspaces,
  $d=\dim_{E}H^{1}(A,\QQ)$; among them are the Weil classes;
\item the Hodge classes of the abelian varieties without isogeny factor $X$;
\item the Hodge classes of the powers of $X$ that are fixed by $\zeta$, among
  them $\mathrm{Det}$ and $r_{1}$.
\end{enumerate}
Then every class of $\cS$ is invariant under $G\cdot A_{3}$ acting through
$V$, and the classes obtained from \textup{(a)}, \textup{(b)} and
\textup{(c)} alone are invariant under $\Sp(V,\psi)$. The Hodge classes in
$\cS(X^{m})$ are the Hodge classes of $X^{m}$ fixed by $\zeta$, and $\cS$
contains no exceptional class of $X\times X$.
\end{proposition}

\begin{proof}
Up to isogeny $A=A_{X}\times A'$ by Poincar\'e reducibility, with $A_{X}$
isogenous to some $X^{k}$ and no isogeny factor $X$ in $A'$, so
$\Hom(X,A')=\Hom(A',X)=0$, $X$ being simple. As $\End^{0}(X)=\QQ$,
$\Hom^{0}(A_{X},X)$ has dimension $k$ and the displayed map is an isomorphism.
A homomorphism $h\colon A\to B$ is, up to isogeny, $h_{X}\times h'$ with
$h_{X}\colon A_{X}\to B_{X}$ and $h'\colon A'\to B'$, and $h^{*}$ acts on
$H^{1}$ as $(1\otimes(f\mapsto f\circ h_{X}))\oplus h'^{*}$, which commutes
with $\GL(V)$; so pull-backs are equivariant, and so are cup and exterior
products, $\GL(V)$ acting on $H^{*}(A\times B)=H^{*}(A)\otimes H^{*}(B)$
diagonally. The orientation class of $A$ spans
${\textstyle\bigwedge^{\mathrm{top}}}H^{1}(A)=(\det V)^{\otimes k}\otimes
(\det\Hom^{0}(A_{X},X))^{\otimes8}\otimes{\textstyle\bigwedge^{\mathrm{top}}}H^{1}(A')$, on
which $\SL(V)$ acts trivially, so Poincar\'e duality and push-forwards are
$\SL(V)$-equivariant. Since $G\cdot A_{3}\subseteq\Sp(V,\psi)\subseteq\SL(V)$,
the classes invariant under either group form a family closed under all the
operations, and it remains to look at the generators.

(a) A divisor class of $A$ is the sum of its K\"unneth components in
${\textstyle\bigwedge^{2}}H^{1}(A_{X})$, $H^{1}(A_{X})\otimes H^{1}(A')$ and
${\textstyle\bigwedge^{2}}H^{1}(A')$, which are Hodge classes. The middle one corresponds to a
homomorphism from $A'$ to the dual of $A_{X}$, which is zero. The first is a
Hodge class of $A_{X}$, so a $G$-invariant of ${\textstyle\bigwedge^{2}}(V\otimes\QQ^{k})=
S^{2}V\otimes{\textstyle\bigwedge^{2}}\QQ^{k}\oplus{\textstyle\bigwedge^{2}}V\otimes S^{2}\QQ^{k}$; as
$(S^{2}V)^{G}=0$ and $({\textstyle\bigwedge^{2}}V)^{G}=\QQ\psi$ (proof of
\Cref{thm:lefschetzclosure}(i)), it lies in
$\psi\otimes S^{2}\QQ^{k}$, which $\Sp(V,\psi)$ fixes. On the third $\GL(V)$
acts trivially. (b) The field $E$ acts on $Q=\Hom^{0}(A_{X},X)$ and on
$H^{1}(A')$, and the $\sigma$-eigenspace of $H^{1}(A,\bar{\QQ})$ is
$V\otimes Q_{\sigma}\oplus H^{1}(A',\bar{\QQ})_{\sigma}$. Its top exterior
power is $(\det V)^{\otimes\dim Q_{\sigma}}\otimes(\det Q_{\sigma})^{\otimes8}
\otimes\det H^{1}(A',\bar{\QQ})_{\sigma}$, on which $\SL(V)$ acts trivially.
(c) For such an abelian variety
$A_{X}=0$, and $\GL(V)$ acts trivially. (d) The Hodge classes of $X^{m}$ are
the invariants of $G$, which acts through $V$, and those fixed by $\zeta$ are
invariant under the group $G\cdot A_{3}$ that $G$ and $\zeta$ generate.

So every class of $\cS$ is $G\cdot A_{3}$-invariant, and those obtained from
(a), (b) and (c) are $\Sp(V,\psi)$-invariant. The Hodge classes of $X^{m}$ in
$\cS$ contain those of (d) and are fixed by $\zeta$, so they are exactly the
Hodge classes fixed by $\zeta$; and an exceptional class of $X\times X$ is not
fixed by $\zeta$ (proof of \Cref{prop:mumfordmotivic}(ii)).
\end{proof}

So no argument that uses only these inputs and these operations reaches the
exceptional classes. The inputs include all Weil classes, whether or not they
are known to be algebraic, and with them Markman's cycles \cite{Mar25a}; the
whole Hodge conjecture for the abelian varieties without isogeny factor $X$;
and the algebraicity of $\mathrm{Det}$ and $r_{1}$. A proof needs a cycle
outside this system: on a product of a power of $X$ with an abelian variety, a
class that is not invariant under $G\cdot A_{3}$ acting through $V$, or a cycle
through a variety that is not an abelian variety, such as the Kuga--Satake
class of $S_{\lambda}$ would provide (\Cref{prop:mumfordkslink}).
\Cref{prop:mumfordformal} extends the second limit of
\Cref{rem:mumfordpowersopen}; it is an obstruction to a method, not to the
target, and we claim no priority for it.

The same invariance, used the other way round, shows that the statement (F2)
of \Cref{cor:fourremain} is not smaller than the conjecture for abelian
varieties. For an abelian variety $Y$ let $R(Y)\subseteq H^{*}(Y,\QQ)$ be the
subring generated by divisor classes and Weil classes of quotients, as in
\Cref{rem:f3collapse}.

\begin{proposition}[Statement (F2) is the conjecture for abelian varieties]
\label{prop:f2isab}
Statement \textup{(F2)} of \Cref{cor:fourremain} implies the Hodge
conjecture for every abelian variety, and is therefore equivalent to it.
\end{proposition}

\begin{proof}
The conjecture for abelian varieties contains (F2). For the converse, recall
from \Cref{rem:f3collapse} that (F2) gives every Hodge class of degree $2q$ on
$Y$ as soon as one Hodge class of degree $2q$ on $Y$ lies outside $R(Y)$. Let
$X$ be as in \Cref{setup:mumford}, and let $\zeta$ act on the cohomology of
every abelian variety through the factor $V$, as in \Cref{prop:mumfordformal}.
Every class of $R(Y)$ lies in $\cS(Y)$, so it is fixed by $\zeta$, for every
$Y$; and an exceptional class $e_{0}\in H^{4}(X\times X,\QQ)$ is not
(\Cref{prop:mumfordformal}). Let $\theta$ be an ample class on $X\times X$,
which $\zeta$ fixes. If $\zeta e_{0}=\lambda e_{0}$, then $\lambda\ne1$ and
$e=e_{0}+\theta^{2}$ has $\zeta e=\lambda e_{0}+\theta^{2}$, which is not a
multiple of $e$; otherwise put $e=e_{0}$. In either case $e$ is a Hodge class
of degree $4$ on $X\times X$ with $\zeta e\notin\bar{\QQ}\,e$.

Let $A$ be an abelian variety and $\beta\ne0$ a Hodge class of degree $2p$ on
$A$, and put $Y=A\times X\times X$. As $\zeta$ acts on
$H^{*}(Y)=H^{*}(A)\otimes H^{*}(X\times X)$ diagonally, a class
$\beta\otimes e$ in $R(Y)$ would satisfy
$\zeta\beta\otimes\zeta e=\beta\otimes e$, which forces $\zeta e\in\bar{\QQ}\,e$.
So the Hodge class $\beta\otimes e$ lies outside $R(Y)$, and (F2) makes it
algebraic. In the same way $e\cup\theta^{4}\notin R(X\times X)$, since
otherwise $(\zeta e-e)\cup\theta^{4}=0$ and the hard Lefschetz theorem would
give $\zeta e=e$; so (F2) makes every Hodge class of degree $12$ on
$X\times X$ algebraic. By the Hodge--Riemann bilinear relations the
intersection pairing between the Hodge classes of degrees $4$ and $12$ on
$X\times X$ is nondegenerate, so there is a Hodge class $e'$ of degree $12$
with $\int_{X\times X}e\cup e'=1$. Then
\[
  \beta=\mathrm{pr}_{A*}\bigl((\beta\otimes e)\cup
  \mathrm{pr}_{X\times X}^{*}e'\bigr),
\]
with $\mathrm{pr}_{A}$ and $\mathrm{pr}_{X\times X}$ the projections of $Y$,
and $\beta$ is algebraic.
\end{proof}

\begin{remark}[The rule set with \Cref{prop:f2isab}]\label{rem:f2rule}
\Cref{thm:closure} is a statement about the rule set of \Cref{setup:rules},
and \Cref{prop:f2isab} is a rule that the set does not contain. Adjoined to
it, the rule changes the fifth minimal set of \Cref{thm:closure}(ii): the
minimal sufficient sets become
\[
  \{\mathrm{F3}\},\quad\{\mathrm{L},\mathrm{M}\},\quad
  \{\mathrm{L},\mathrm{F3}'\},\quad\{\mathrm{V},\mathrm{F3}'\},\quad
  \{\mathrm{F2},\mathrm{F3}'\},
\]
and $\mathrm{P2}_{\mathrm{split}}$ lies in none of them. Indeed the new rule
has the premise (F2) and the conclusion $\mathbf{AB}$, the conjecture for
abelian varieties, so the sets $C_{1}$, $C_{2}$ and $C_{4}$ of the proof of
\Cref{thm:closure} stay closed, and a sufficient set meets $H_{1}$, $H_{2}$
and $H_{4}$. The sets meeting these three that are minimal with this property
are the five above, by the argument of that proof without $H_{3}$, and each is
sufficient: $\{\mathrm{F2},\mathrm{F3}'\}$ gives $\mathbf{AB}$ by the new
rule and then $\mathbf{HC}$ by \Cref{prop:f3prime}(i). So the results of this
paper on Weil classes do not shorten a route to the conjecture. What they do
is reduce one case of it, the algebraicity of the Weil classes of every CM
field, to the single statement $\mathrm{P2}_{\mathrm{split}}$, by a route that
does not pass through (F2). Like (F3), which is the whole conjecture, (F2) is
the whole conjecture for abelian varieties, carried by the classes that are
hardest to reach; the exceptional classes of $X\times X$ are its smallest
instance, and the proof above shows that the classes $\beta\otimes e$ built
from them already carry all of it.
\end{remark}
